\question{10.m1}{Voor een kansvariabele die een theoretische chikwadraatsverdeling met vijf vrijheidsgraden volgt, geldt dat $\ldots$}
\begin{enumerate}[label=(\alph*)]
    \item deze symmetrisch rondom $0$ ligt.
    \item deze symmetrisch rondom $5$ ligt.
    \item de kansen hiervoor kunnen worden gevonden als het kwadraat wordt genomen van een normaal verdeelde variabele met $\mu=5$.
    \item deze kansvariabele waarden kan aannemen die groter zijn dan $15$.
\end{enumerate}
\answer{
    De uitkomstenruimte van een kansvariabele met een chikwadraatsverdeling is altijd gelijk aan alle reële getallen groter dan of gelijk aan nul,
    dus ook waarden die groter zijn dan $15$.
    Het juiste antwoord is dus (d).
}
